\begin{table}[htbp]\centering
\caption{全部背景空隙明细（整段 RMS vs 安静帧）}\label{tab:nm_bg_detail}
{\footnotesize\begin{threeparttable}
\begin{tabular}{lccccccc}\toprule
片段 & 起点 & 整段旧 & 整段新 & 安静帧旧 & 安静帧新 & 峰值 & 判定 \\
\midrule
test2 & 0.3 & $-$5.6 & $-$6.7 & $-$6.2 & $-$6.9 & $-$56 & 底噪未超阈值 \\
test2 & 3.5 & $-$4.9 & $-$4.0 & $-$6.0 & $-$7.5 & $-$55 & 底噪未超阈值 \\
test4 & 4.3 & $+$3.8 & $-$0.5 & $+$1.9 & $-$4.4 & $-$34 & 含突发声（非整段底噪） \\
test12 & 4.3 & $-$0.5 & $+$0.4 & $-$0.6 & $-$2.5 & $-$60 & 底噪未超阈值 \\
test13 & 5.8 & $+$0.5 & $+$0.2 & $+$0.5 & $+$0.1 & $-$39 & 含突发声（非整段底噪） \\
test14 & 14.5 & $+$13.3 & $+$14.0 & $+$6.4 & $+$0.8 & $-$39 & 含突发声（非整段底噪） \\
\bottomrule\end{tabular}
\begin{tablenotes}\footnotesize\item 含突发声：段内短时峰值 $>-55$ dBFS，整段增益不能代表底噪。\end{tablenotes}\end{threeparttable}}
\end{table}
